\chapter{Motion Capture}
\label{section_mocap}

\section{Installation du package \textit{vrpn\_mocap}}

La section suivante décrit comment récupérer, sur un topic ROS2, la position des \textit{Rigid Bodies} détectés par le système de Motion Capture Motive.

Commencer par créer un nouvel espace de travail ROS2. Ouvrir un terminal et entrer les commandes suivantes pour créer le répertoire \textit{mocap\_ws} ainsi que son sous-répertoire \textit{src} :

\begin{lstlisting}[style=commandstyle]
mkdir -p ~/mocap_ws/src
cd ~/mocap_ws
\end{lstlisting}

\noindent Cloner ensuite le dépôt GitHub du package \textit{vrpn\_mocap} dans le répertoire \textit{src} de l’espace de travail :

\begin{lstlisting}[style=commandstyle]
cd ~/mocap_ws/src
git clone https://github.com/alvinsunyixiao/vrpn_mocap.git
\end{lstlisting}

\noindent Une fois le dépôt cloné, installer les dépendances nécessaires à la compilation en exécutant la commande suivante :

\begin{lstlisting}[style=commandstyle]
rosdep install --from-paths src -y --ignore-src
\end{lstlisting}

\noindent Procéder ensuite à la compilation du package :

\begin{lstlisting}[style=commandstyle]
colcon build
\end{lstlisting}

\noindent Enfin, sourcer l’espace de travail :

\begin{lstlisting}[style=commandstyle]
source install/setup.bash
\end{lstlisting}

Pour éviter d’avoir à sourcer l’environnement à chaque nouvelle session, ajouter cette ligne au fichier \textit{.bashrc} :

\begin{lstlisting}[style=commandstyle]
echo "source ~/mocap_ws/install/setup.bash" >> ~/.bashrc
\end{lstlisting}

Il est possible de personnaliser le comportement du package en modifiant le fichier \textit{client.launch.yaml}. Les principaux paramètres configurables sont les suivants :

\begin{itemize}
    \item \textit{server} : adresse IP ou nom de domaine du serveur VRPN (par défaut : \textit{localhost}).
    \item \textit{port} : port utilisé par le serveur VRPN (par défaut : \textit{3883}).
    \item \textit{frame\_id} : nom de la frame de référence fixe (par défaut : \textit{world}).
    \item \textit{update\_freq} : fréquence de publication des données de capture de mouvement (par défaut : \textit{100.0 Hz}).
\end{itemize}

Pour tester le package, utiliser la commande suivante (modifier l'adresse IP et le port pour correspondre au système de Motion Capture) : 
\begin{lstlisting}[style=commandstyle]
ros2 launch vrpn_mocap client.launch.yaml server:=192.168.0.103 port:=3883
\end{lstlisting}

Il est possible de modifier le fichier \textit{client.launch.yaml}, afin de ne pas avoir à spécifier l'IP du serveur et le port dans la commande. Pour cela, modifier le fichier \textit{client.launch.yaml} se trouvant dans \textit{/mocap/src/vrpn\_mocap/launch} pour qu'il corresponde à : 
\begin{lstlisting}[style=commandstyle]
launch:

- arg:
    name: "server"
    default: "192.168.0.103"
- arg:
    name: "port"
    default: "3883"

- node:
    pkg: "vrpn_mocap"
    namespace: "vrpn_mocap"
    exec: "client_node"
    name: "vrpn_mocap_client_node"
    param:
    -
      from: "$(find-pkg-share vrpn_mocap)/config/client.yaml"
    -
      name: "server"
      value: "$(var server)"
    -
      name: "port"
      value: "$(var port)"

\end{lstlisting}

\noindent La nouvelle commande à utiliser pour lancer le noeud sera alors :
\begin{lstlisting}[style=commandstyle]
ros2 launch vrpn_mocap client.launch.yaml
\end{lstlisting}
\label{lancement_mocap}

\section{Quality of Service (QoS) sur ROS2}

Les paramètres de \textit{Quality of Service (QoS)} sur ROS2 permettent de contrôler la manière dont les données sont échangées entre les nœuds.
Par exemple, le paramètre \textit{reliability} spécifie la garantie de livraison des messages entre les publishers et les subscribers. Deux options sont disponibles pour ce paramètre :

\begin{itemize}
    \item \textbf{Best Effort} : tente de livrer les messages, mais certains peuvent être perdus si le réseau est instable. Ce mode est analogue au protocole \textbf{UDP}, qui privilégie la rapidité d'envoi des paquets sans garantie de réception. 
    \item \textbf{Reliable} : garantit la livraison des messages, quitte à effectuer plusieurs tentatives en cas d’échec. Ce comportement est similaire à \textbf{TCP}, qui assure la fiabilité de la transmission en retransmettant les données en cas de perte.
\end{itemize}


\noindent Cependant, ces deux modes ne sont pas toujours compatibles entre eux. Par exemple, un \textit{subscriber} configuré en \textit{Reliable} ne pourra pas recevoir les messages d’un \textit{publisher} en mode \textit{Best Effort}, car il s’attend à une garantie stricte de recevoir des données. En ROS2, les \textit{publisher} et \textit{subscriber} sont initialement réglés en \textit{Reliable}, et ne peuvent donc pas afficher les messages publiés par \textit{vrpn\_mocap} qui sont en \textit{Best Effort}. Pour régler ce problème, il faut, si possible, modifier les configurations QoS (voir Annexe \ref{annexe:qos}). 

\noindent Ci-dessous un tableau récapitulatif des compatibilités du paramètre \textit{Reliable :}
\begin{table}[h]
    \centering
    \begin{tabular}{|c|c|c|}
        \hline
        \textbf{Publisher} & \textbf{Subscriber} & \textbf{Compatible} \\
        \hline
        Best Effort & Best Effort & Oui \\
        Best Effort & Reliable & Non \\
        Reliable & Best Effort & Oui \\
        Reliable & Reliable & Oui \\
        \hline
    \end{tabular}
    \caption{Compatibilité du paramètre \textit{Reliable }}
    \label{tab:qos_compatibility}
\end{table}

Si modifier les paramètres n'est pas possible (comme par exemple sur rqt), il est nécessaire de d'ajouter la ligne ci-dessous dans le fichier \textit{client.yaml}, situé dans \textit{mocap\_ws/src/vrpn\_mocap/config}, afin de forcer le paramètre \textit{Reliability} de \textit{vrpn\_mocap} en \textit{Reliable} :

\begin{lstlisting}[style=commandstyle]
sensor_data_qos: false
\end{lstlisting}

\noindent Le fichier \textit{client.yaml} final doit ressembler à ceci :

\begin{lstlisting}[style=commandstyle]
/vrpn_mocap/vrpn_mocap_client_node:
  ros__parameters:
    server: localhost
    port: 3883
    frame_id: world
    update_freq: 100.
    refresh_freq: 1.
    sensor_data_qos: false
    multi_sensor: false
    use_vrpn_timestamps: false
\end{lstlisting}

\noindent 
Puis, \textit{build} à nouveau pour être sur de prendre en compte les modifications.
\begin{lstlisting}[style=commandstyle]
colcon build
source install/setup.bash
\end{lstlisting}


Une alternative, dans le cas où l'utilisation du mode \textit{Best Effort} est impérative pour la publication, consiste à créer un nœud intermédiaire contenant un \textit{subscriber} configuré en \textit{Best Effort}. Ce dernier souscrira aux topics fournis par \textit{vrpn\_mocap} et republiera les messages à l’aide d’un \textit{publisher} configuré en \textit{Reliable}, sur de nouveaux topicx. Cela permet ainsi d'assurer la compatibilité avec les outils nécessitant une QoS fiable, tout en conservant les topics réglés en \textit{Best Efforts}.

Des informations supplémentaires sur les QoS sont disponibles ici : \newline
\url{https://docs.ros.org/en/jazzy/Concepts/Intermediate/About-Quality-of-Service-Settings.html}
